\section{Related work}\label{sec:related}

\textbf{Machine learning potentials}:
Machine learning has a long history in learning potentials and force fields for molecular simulations. Starting with \citet{halgrenMerckMolecularForce1996} where empirical force fields were introduced with the goal of fitting functions to QM calculations to avoid performing the expensive QM calculations many times.
Later, kernel methods~\citep{behlerAtomcenteredSymmetryFunctions2011,bartokRepresentingChemicalEnvironments2013,christensenFCHLRevisitedFaster2020} based on sophisticated feature constructions of molecular neighborhoods took over the field.
These kernel methods were fast to evaluate and capable of accurately capturing molecular interactions.
By focusing on locality and freely learnable featurizations GNNs~\citep{schuttSchNetDeepLearning2018,unkePhysNetNeuralNetwork2019} achieved accurate reproductions of molecular force fields.
Due to the implementation of physical symmetries to the euclidean group $\set E(3)$ and the permutational group $S_n$, GNNs presented sample efficient ways of learning from molecular data.
Recent advances provided universal models, i.e., models that can model any function on geometry point clouds, accomplished by including dihedral angles~\citep{gasteigerGemNetUniversalDirectional2021,liuSphericalMessagePassing2021} or by the introduction of SO(3) equivariant models~\citep{thomasTensorFieldNetworks2018,batznerEquivariantGraphNeural2022}.
Other developments focus on closing the gap between QM calculations and ML potentials, e.g., by including prior physical knowledge~\citep{unkeSpookyNetLearningForce2021}, by providing efficient QM calculations as additional input~\citep{qiaoOrbNetDeepLearning2020,qiaoUNiTEUnitaryNbody2021}, by implementing DFT functionals with GNNs~\citep{snyderFindingDensityFunctionals2012,kirkpatrickPushingFrontiersDensity2021}, or by learning potentials directly from first principle~\citep{gaoAbInitioPotentialEnergy2022,gaoGeneralizing2023,gaoSamplingfreeInferenceAbInitio2022}.

\textbf{Uncertainty in machine learning.} There are different approaches for uncertainty estimation (UE) in deep learning. The interested reader is referred to \citet{gawlikowski2021survey} for a detailed survey. UE methods can be broadly categorized into two families of methods: \emph{sampling-based} and \emph{sampling-free} methods. 

Sampling-based methods generally evaluate the uncertainty by estimating statistics over multiple different predictions. Generally, when the prediction is very similar of over the different samples, this indicates the low-uncertainty of the model. Ensembles \citep{bayesian-classifier-combination,ensembles, dynamic-bayesian-combination-classifiers,batch-ensembles,hyper-ensembles}, often refered to as the gold-standard, are a collection of individual models. They estimate the uncertainty by the variation in the prediction of the different ensemble members. Dropout methods \cite{gal2016dropout}, approximate this technique by following the argument that dropout let the model learn different subnetworks during training. Thus each random initialization might be an approximation of an ensemble member. Other notable techniques are bayesian neural networks \citep{bayesian-networks, scalable-laplace-bnn, simple-baseline-uncertainty}, where the model parameters are a distribution instead of a single value. Hence, one can sample parameters of the model multiple times and use these to infer a data sample resulting in multiple predictions from which we can calculate the variance again. 

Sampling-free generally requires to evaluate the uncertainty in a single forward-pass methods. Evidential methods propose to parametrize conjugate prior distributions \citep{survey_evidential_uncertainty,dbu-robustness,max_gap_id_ood,uncertainty-generative-classifier,multifaceted_uncertainty,graph_posterior, lightweight-prob-net}. Deep Kernel Learning aim at learning Gaussian processes in latent space with random feature projections or learned inducing points \citep{uncertainty-distance-awareness, due, duq, uceloss}. Calibration models are methods which aim at predicting confidence which are good approximations of the true probability of correctness \citep{accurate-uncertainties-deep-learning-regression, confidence-aware-learning, individual-calibration, distribution-calibration-regression, intra-order-preserving}. Other models propose to propagate uncertainty accross layers and model uncertainty at the weight or activation levels \citep{natural-parameter-network, sampling-free-variance-propagation, feed-forward-propagation, lightweight-prob-net, probabilistic-backprop-scalable-bnn}. Most of these models do not predict uncertainty on molecular properties. Only \citet{soleimany2021evidential} leverage uncertainty estimation of evidential deep learning methods for guided molecular property prediction and show its effectiveness. We use their work as a further baseline, to which we compare our results. Another interesting approach is the Graph Mixture Density Network \cite{errica2021graph}, which outputs a mixture of distributions on the predicted target but does not focus on molecules.


